\subsection*{Format du rendu --- en-tête obligatoire}

Tout fichier rendu dans ce module doit commencer par l'en-tête
ci-dessous, \textbf{adapté à la syntaxe de commentaire du langage}
utilisé.

\subsubsection*{Modèle générique}

\begin{minted}{c}
/**
 * Module: <code module>
 * Exercise: <numéro exo>
 * File: <nom du fichier>
 * Author: <votre login JeuneCrack>
 *
 * Description: <une phrase, pas de roman>
 *
 * Return: <voir ci-dessous>
 */
\end{minted}

\subsubsection*{Login}

Votre identifiant \textbf{JeuneCrack} (ex: \texttt{tedsig42}).
Pour un rendu en binôme : \texttt{login1/login2}.
Ajoutez un handle Matrix après le login si vous le souhaitez
(ex: \texttt{tedsig42 @tedsig42}).

\subsubsection*{Return --- selon le type de fichier}

\begin{tabularx}{\linewidth}{lX}
  \toprule
  \textbf{Type de fichier} & \textbf{Valeur de Return} \\
  \midrule
  C (\texttt{.c}) &
  \texttt{return (0);} (succès),
  \texttt{return (-1);} (erreur),
  \texttt{return (valeur);} (résultat calculé) \\

  Python (\texttt{.py}) &
  \texttt{0} (succès),
  \texttt{1} (erreur),
  \texttt{N/A} (analyse, script sans exit) \\

  Shell (\texttt{.sh}) &
  \texttt{exit 0} (succès),
  \texttt{exit 1} (erreur) \\

  Markdown (\texttt{.md}) &
  \texttt{N/A} --- les fichiers d'analyse n'ont pas de valeur de retour \\

  Exploit C/Python &
  \texttt{shell} (shell obtenu),
  \texttt{flag: <flag>} (flag capturé),
  \texttt{proof: <hash>} (preuve de concept) \\

  Binaire (\texttt{.bin}) &
  \texttt{N/A} --- pas d'en-tête pour les fichiers binaires \\
  \bottomrule
\end{tabularx}

\subsubsection*{Exemples par langage}

\textbf{C (\texttt{.c})} :
\begin{minted}{c}
/**
 * Module: F-EXP-001
 * Exercise: ex02
 * File: exploit.c
 * Author: tedsig42
 *
 * Description: Buffer overflow with ret2libc
 *
 * Return: 0
 */
\end{minted}

\textbf{Python (\texttt{.py})} :
\begin{minted}{python}
# Module: F-WEB-001
# Exercise: ex05
# File: ssrf-exploit.py
# Author: tedsig42
#
# Description: Blind SSRF bypass via redirect
#
# Return: 0
\end{minted}

\textbf{Shell (\texttt{.sh})} :
\begin{minted}{bash}
# Module: F-SUP-001
# Exercise: ex03
# File: dependency-confusion.sh
# Author: tedsig42
#
# Description: Automated dependency confusion test
#
# Return: exit 0
\end{minted}

\textbf{Markdown (\texttt{.md})} --- analyse, rapport :
\begin{Verbatim}[fontsize=\normalsize]
<!--
Module: M-SRC-001
Exercise: ex04
File: audit-report.md
Author: tedsig42

Description: Code audit of vulnerable-httpd

Return: N/A
-->
\end{Verbatim}
